\section{Threat Model and Security Analysis}
\label{s:threat-model}
Our design assumes a decentralized network of organizations where a subset of nodes may be faulty or malicious. We describe the trust assumptions for each entity and analyze the security implications when these assumptions are violated.

\textit{Peer Nodes.} We assume a crash-fault tolerant (CFT) model for authorized peers. A malicious peer may attempt to (a)~leak private data it is authorized to see, or (b)~refuse to disseminate data, but a single compromised peer cannot manipulate the ledger state because the integrity of the public ledger is protected by the endorsement policy, which requires signatures from multiple organizations before a transaction is considered valid. Data manipulation is only possible if a majority of the endorsing organizations collude. \textit{Correctness} is preserved as long as the endorsement policy holds---that is, a sufficient number of endorsing organizations remain honest. If a majority of the endorsing organizations collude, they can endorse fraudulent transactions with manipulated hashes, compromising both the private data and the on-chain evidence. Data \textit{confidentiality} within a collection is bounded by the weakest member---a compromised member can exfiltrate any private data it is authorized to access. \textbf{This is an inherent limitation of any access-control-based (as opposed to cryptographic) privacy mechanism}: the system ensures that unauthorized parties never receive private data, but it cannot prevent an authorized party from misusing data it legitimately possesses. Organizations mitigate this risk through legal agreements, hardware security modules, and operational controls.

\textit{Ordering Service.} The ordering service is a decentralized cluster operated by consortium organizations---potentially the same organizations that run peers. Its role is strictly limited to establishing the total order of transactions; it never sees private data payloads, as only public and hashed read-write sets are submitted for ordering. Standard BFT security assumptions apply: as long as fewer than one-third of the ordering nodes are Byzantine, the ordering service guarantees liveness and correct ordering.

\textit{Clients.} Clients may be malicious, attempting to submit inconsistent private data to different endorsers or forge simulation results. If a client sends different private data to different endorsers, each endorser will compute a different hash of the private data and include it in its signed endorsement response. Since the endorsement policy requires matching responses from multiple endorsers, the inconsistent hashes will cause the endorsements to be incompatible, and the transaction will be rejected during validation. A malicious client may also simulate many transactions but never submit them to the ordering service, leading to the accumulation of orphan write-sets in the transient store. To mitigate this, we introduce a retention window $\Delta$: the transient store automatically purges any entries where the current ledger height exceeds the simulation height by more than $\Delta$. Receiving peers also reject disseminated write-sets outside this window. A second attack vector involves fabricated simulations lacking sufficient endorsements; to counter this, peers perform endorsement policy validation \emph{before} attempting to pull any missing pre-images, skipping invalid transactions entirely.

\textit{Network.} We assume a partially synchronous network where partitions are possible but eventually resolve, allowing for asynchronous data reconciliation. During a partition, private data dissemination may be delayed but the system does not compromise correctness---peers that cannot obtain private data before commit track it as missing and reconcile asynchronously once connectivity is restored.

